% English translation of chapter1.tex lines 1141--1229.
% Source: Methods of Algebra, Volume 2, commit
% 9a5803ff77dd3257484cb177f851a73770a59dd3 (CC BY 4.0).
% stable-unit-id: o014.aljabr2.chapter1.kan-extensions-as-limits

% segment-id: o014.li2.chapter1.kan-extensions-as-limits.h001
\section{Constructing Kan Extensions by Limits}\label{sec:Kan-as-limit}
\hypertarget{o014.aljabr2.chapter1.kan-extensions-as-limits}{ }

% segment-id: o014.li2.chapter1.kan-extensions-as-limits.p001
As in \S\ref{sec:Kan-extension}, throughout this section we continue to
consider functors $K: \mathcal{C} \to \mathcal{D}$ and
$F: \mathcal{C} \to \mathcal{E}$, with the aim of studying the existence of
$\Lan_K F$ and $\Ran_K F$.

% segment-id: o014.li2.chapter1.kan-extensions-as-limits.p002
Recall the discussion at the beginning of \S\ref{sec:Kan-extension}. The
motivation for Kan extensions is to seek a functor
$G: \mathcal{D} \to \mathcal{E}$ such that $F \simeq GK$, or at least to find
the best approximation to one. How should $Gd$ be defined for
$d \in \Obj(\mathcal{D})$? The only clue is that when $d=Kc$, the object $Gd$
ought to be $Fc$, up to isomorphism. This suggests approximating $Gd$ by a
colimit $\varinjlim$ (or a limit $\varprojlim$) of the objects $Fc$, indexed by
	all $c \in \Obj(\mathcal{C})$ and morphisms $Kc \to d$ (or $d \to Kc$).
	These two constructions have dual universal properties. We now make the
	construction precise.

% segment-id: o014.li2.chapter1.kan-extensions-as-limits.p003
Given $d \in \Obj(\mathcal{D})$, define the categories $(K/d)$ and $(d/K)$ as
in Definition \ref{def:comma-category}, together with the projection functors
% segment-id: o014.li2.chapter1.kan-extensions-as-limits.q001
\[ \Pi_{/d}: (K/d) \to \mathcal{C}, \quad \Pi_{d/}: (d/K) \to \mathcal{C}. \]
Our focus will be on the composite functors
$F\Pi_{/d}: (K/d) \to \mathcal{E}$ and
$F\Pi_{d/}: (d/K) \to \mathcal{E}$, which send the objects
$(c,Kc \to d)$ and $(c,d \to Kc)$, respectively, to $Fc$. Assuming that the
limits under discussion exist, we first record a few observations.

% segment-id: o014.li2.chapter1.kan-extensions-as-limits.l001
\begin{itemize}
	% segment-id: o014.li2.chapter1.kan-extensions-as-limits.i001
	\item Every morphism $d \to d'$ induces a canonical morphism
	$\varinjlim F\Pi_{/d} \to \varinjlim F\Pi_{/d'}$; it is characterized by
	requiring the diagram
	% segment-id: o014.li2.chapter1.kan-extensions-as-limits.q002
	\begin{equation}\label{eqn:Kan-limit-aux-0}
	% segment-id: o014.li2.chapter1.kan-extensions-as-limits.g001
	\begin{tikzcd}
		Fc \arrow[equal, r] \arrow[d] & Fc \arrow[d] \\
		\varinjlim F \Pi_{/d} \arrow[r] & \varinjlim F \Pi_{/d'}
	\end{tikzcd}\end{equation}
	to commute for every $(c,Kc \to d) \in \Obj((K/d))$ (which induces an
	object $(c,Kc \to d') \in \Obj((K/d'))$), where the vertical arrows are
	the structure morphisms of the colimits. This is an instance of the
	functoriality of colimits; see \cite[Lemma 2.7.4]{Li1}\footnote{Translator's
	note: the source gives the locator as Lemma~2.7,4; it has been corrected to
	Lemma~2.7.4 here (O014-C016).} and its proof.
	% segment-id: o014.li2.chapter1.kan-extensions-as-limits.i002
	\item Considering
	$(c,Kc \xrightarrow{\identity} Kc) \in \Obj(K/Kc)$ gives a canonical
	morphism $\eta_c: Fc \to \varinjlim F\Pi_{/Kc}$.
	% segment-id: o014.li2.chapter1.kan-extensions-as-limits.i003
	\item Every morphism $c \to c'$ induces a canonical morphism
	$\varinjlim F\Pi_{/Kc} \to \varinjlim F\Pi_{/Kc'}$ that makes the following
	diagram commute:
	% segment-id: o014.li2.chapter1.kan-extensions-as-limits.g002
	\[\begin{tikzcd}
		Fc \arrow[r] \arrow[d, "{\eta_c}"'] & Fc' \arrow[d, "{\eta_{c'}}"] \\
		\varinjlim F \Pi_{/Kc} \arrow[r] & \varinjlim F \Pi_{/Kc'} .
	\end{tikzcd}\]
\end{itemize}

% segment-id: o014.li2.chapter1.kan-extensions-as-limits.p004
By duality, $\varprojlim F\Pi_{d/}$ has the corresponding properties,
including a canonical morphism
$\varepsilon_c: \varprojlim F\Pi_{Kc/} \to Fc$. The next theorem is based on
this functoriality.

% segment-id: o014.li2.chapter1.kan-extensions-as-limits.th001
\begin{theorem}\label{prop:Kan-limit}
	Consider functors $K: \mathcal{C} \to \mathcal{D}$ and
	$F: \mathcal{C} \to \mathcal{E}$.
	% segment-id: o014.li2.chapter1.kan-extensions-as-limits.l002
	\begin{enumerate}[(i)]
		% segment-id: o014.li2.chapter1.kan-extensions-as-limits.i004
		\item Suppose that for every $d \in \Obj(\mathcal{D})$, the colimit
		$\varinjlim(F\Pi_{/d})$ exists in $\mathcal{E}$. Then
		$(\Lan_K F)(d) := \varinjlim(F\Pi_{/d})$, together with the
		functoriality given by \eqref{eqn:Kan-limit-aux-0}, determines the left
		Kan extension $\Lan_K F: \mathcal{D} \to \mathcal{E}$. The associated
		natural transformation $\eta: F \to (\Lan_K F)\circ K$ comes from the
		family $(\eta_c)_{c\in\Obj(\mathcal{C})}$ discussed above.
		% segment-id: o014.li2.chapter1.kan-extensions-as-limits.i005
		\item Suppose that for every $d \in \Obj(\mathcal{D})$, the limit
		$\varprojlim(F\Pi_{d/})$ exists in $\mathcal{E}$. Then
		$(\Ran_K F)(d) := \varprojlim(F\Pi_{d/})$, together with the dual of
		\eqref{eqn:Kan-limit-aux-0}, determines the right Kan extension
		$\Ran_K F: \mathcal{D} \to \mathcal{E}$. The associated natural
		transformation $\varepsilon: (\Ran_K F)\circ K \to F$ comes from the
		family $(\varepsilon_c)_{c\in\Obj(\mathcal{C})}$ discussed above.
		% segment-id: o014.li2.chapter1.kan-extensions-as-limits.i006
		\item If $K: \mathcal{C} \to \mathcal{D}$ is fully faithful, then the
		$\eta$ constructed in (i) (or the $\varepsilon$ constructed in (ii)) is
		an isomorphism.
	\end{enumerate}
\end{theorem}
% segment-id: o014.li2.chapter1.kan-extensions-as-limits.pf001
\begin{proof}
	% segment-id: o014.li2.chapter1.kan-extensions-as-limits.p005
	By duality, we treat only the case of $\Lan_K F$. For (i), we verify the
	universal property in Definition \ref{def:Kan-extension} step by step.

	% segment-id: o014.li2.chapter1.kan-extensions-as-limits.p006
	Let $L: \mathcal{D} \to \mathcal{E}$ be a functor and let
	$\xi: F \to LK$. For every $d \in \Obj(\mathcal{D})$, every object
	$(c,Kc \to d)$ of $(K/d)$, and every morphism in $(K/d)$ determined by
	$f:c\to c'$, namely $(c,Kc\to d)\to(c',Kc'\to d)$, there is a commutative
	diagram
	% segment-id: o014.li2.chapter1.kan-extensions-as-limits.g003
	\[\begin{tikzcd}
		Fc \arrow[r, "\xi_c"] \arrow[d, "Ff"'] & LKc \arrow[d, "LKf"'] \arrow[r] & Ld . \\
		Fc' \arrow[r, "{\xi_{c'}}"'] & LKc' \arrow[ru] &
	\end{tikzcd}\]
	The universal property of $\varinjlim$ therefore gives
	$\chi_d:(\Lan_K F)(d)\to Ld$, characterized by requiring the diagram
	% segment-id: o014.li2.chapter1.kan-extensions-as-limits.q003
	\begin{equation}\label{eqn:Kan-limit-aux-1}
	% segment-id: o014.li2.chapter1.kan-extensions-as-limits.g004
	\begin{tikzcd}
		Fc \arrow[r] \arrow[d, "\xi_c"'] & (\Lan_K F)(d) \arrow[d, "\chi_d"] \\
		LKc \arrow[r, "{L(Kc \to d)}"'] & Ld
	\end{tikzcd}\end{equation}
	to commute for all $(c,Kc\to d)\in\Obj((K/d))$; the arrow in the first row
	is the structure morphism of the colimit.

	% segment-id: o014.li2.chapter1.kan-extensions-as-limits.p007
	We now prove that $(\chi_d)_{d\in\Obj(\mathcal{D})}$ defines a natural
	transformation $\chi:\Lan_K F\to L$. Given $d\to d'$, take
	$(c,Kc\to d)$ and its image $(c,Kc\to d')$. Comparing
	\eqref{eqn:Kan-limit-aux-0} with \eqref{eqn:Kan-limit-aux-1}, the question
	reduces to proving that the diagram
	% segment-id: o014.li2.chapter1.kan-extensions-as-limits.g005
	\[\begin{tikzcd}[row sep=tiny]
		& & Ld \arrow[dd, "{L(d \to d')}"] \\
		Fc \arrow[r, "\xi_c"] & LKc \arrow[ru, "{L(Kc \to d)}"] \arrow[rd, "{L(Kc \to d')}"'] & \\
		& & Ld'
	\end{tikzcd}\]
	commutes, which is immediate.

	% segment-id: o014.li2.chapter1.kan-extensions-as-limits.p008
	The next step is to verify
	% segment-id: o014.li2.chapter1.kan-extensions-as-limits.q004
	\begin{equation}\label{eqn:Kan-limit-aux-2}
		\forall c \in \Obj(\mathcal{C}), \quad \xi_c = \chi_{Kc} \eta_c: Fc \to LKc.
	\end{equation}
	Since $\eta_c$ is the canonical morphism from
	$Fc=F\Pi_{/Kc}(c,Kc\xrightarrow{\identity}Kc)$ to
	$(\Lan_K F)(Kc)$, equation \eqref{eqn:Kan-limit-aux-2} is simply a special
	case of \eqref{eqn:Kan-limit-aux-1}.

	% segment-id: o014.li2.chapter1.kan-extensions-as-limits.p009
	Finally, we show that a natural transformation $\chi:\Lan_K F\to L$
	satisfying \eqref{eqn:Kan-limit-aux-2} is unique. Fix $d$ and consider any
	object $(c,Kc\to d)$ of $(K/d)$ and the associated diagram
	% segment-id: o014.li2.chapter1.kan-extensions-as-limits.g006
	\[\begin{tikzcd}
		& Fc \arrow[ld, "{\eta_c}"'] \arrow[rd] & \\
		(\Lan_K F)(Kc) \arrow[rr, "{(\Lan_K F)(Kc \to d)}"'] \arrow[d, "{\chi_{Kc}}"'] & & (\Lan_K F)(d) \arrow[d, "\chi_d"] \\
 		LKc \arrow[rr, "{L(Kc \to d)}"'] & & Ld
	\end{tikzcd}\]
	where $Fc\to(\Lan_K F)(d)$ is the morphism appearing in
	\eqref{eqn:Kan-limit-aux-1}. The square commutes by the naturality of
	$\chi$, while the triangle commutes by \eqref{eqn:Kan-limit-aux-0}. Since
	$\chi_{Kc}\eta_c=\xi_c$, the composite
	$Fc\to(\Lan_K F)(d)\xrightarrow{\chi_d}Ld$ equals
	$L(Kc\to d)\xi_c$. As $(c,Kc\to d)$ varies, this family of equalities
	uniquely determines $\chi_d$.

	% segment-id: o014.li2.chapter1.kan-extensions-as-limits.p010
	Thus $\left(\Lan_K F,\eta\right)$ is indeed a left Kan extension.

	% segment-id: o014.li2.chapter1.kan-extensions-as-limits.p011
	Now consider (iii), and fix $c\in\Obj(\mathcal{C})$. Since $K$ is fully
	faithful, it is easy to see that $(K/Kc)$ has the terminal object
	$(c,Kc\xrightarrow{\identity}Kc)$. The morphism
	$\eta_c:Fc\to\varinjlim F\Pi_{/Kc}$ comes from this terminal object and is
	therefore an isomorphism.
\end{proof}

% segment-id: o014.li2.chapter1.kan-extensions-as-limits.r001
\begin{remark}
	If $\mathcal{C}$ is a small category, then $(K/d)$ and $(d/K)$ are also
	small categories for every $d\in\Obj(\mathcal{D})$. Consequently, when
	$\mathcal{E}$ is cocomplete (or complete), Theorem \ref{prop:Kan-limit} and
	Proposition \ref{prop:Kan-uniqueness} show that
	$K^*: \mathcal{E}^{\mathcal{D}} \to \mathcal{E}^{\mathcal{C}}$ has a left
	adjoint $\Lan_K$ (or a right adjoint $\Ran_K$).
\end{remark}

% segment-id: o014.li2.chapter1.kan-extensions-as-limits.e001
\begin{example}[Inverse image of presheaves]
	\index{presheaf}
	Take $\mathcal{E}:=\cate{Set}$, which is both complete and cocomplete. Let
	$f:X\to Y$ be a continuous map of topological spaces. Let
	$\mathcal{C}:=\cate{Open}_Y$ be the category whose objects are the open
	subsets of $Y$ and whose morphisms are inclusions of open subsets; similarly,
	let $\mathcal{D}:=\cate{Open}_X$. Now define a functor
	$\cate{Open}_Y\to\cate{Open}_X$ that sends an open subset $V\subset Y$ to
	$f^{-1}V$, and regard it below as a functor
	$K:\cate{Open}_Y^{\opp}\to\cate{Open}_X^{\opp}$. Readers familiar with
	sheaf theory will recognize that the objects of
	$\cate{Open}_Y^\wedge:=\cate{Set}^{\cate{Open}_Y^{\opp}}$ are, by
	definition, precisely the \emph{presheaves} on $Y$, while
	$K^*:\cate{Open}_X^\wedge\to\cate{Open}_Y^\wedge$ is the direct-image
	functor between the presheaf categories: it sends a presheaf $\mathcal{F}$
	on $X$ to the presheaf on $Y$ given by
	$V\mapsto\mathcal{F}(f^{-1}V)$. In this situation, the left adjoint of the
	direct image is $\Lan_K$, given by Theorem \ref{prop:Kan-limit}; it sends a
	presheaf $\mathcal{G}$ on $Y$ to
	% segment-id: o014.li2.chapter1.kan-extensions-as-limits.q005
	\begin{gather*}
		U \mapsto \varinjlim_{\substack{V \subset Y: \text{open subset} \\ \text{such that}\; U \subset f^{-1} V }} \mathcal{G}(V), \quad U \subset X: \text{open subset}, \\
		U \subset f^{-1}V \;\xleftrightarrow{\text{corresponds to}}\; \text{a morphism}\; K(V) \to U \;\text{in}\; \cate{Open}_X^{\opp}.
	\end{gather*}
	Unsurprisingly, this is precisely the \emph{inverse image} of $\mathcal{G}$
	in sheaf theory.
\end{example}
